% Sec. Related Work
\section{Related Work}
\label{sec:related}

A team of robots and fixed sensors that shares perception over Wi-Fi
is an IoT system: heterogeneous nodes sense, compute, and exchange
data over one shared, contended network, and the application succeeds
only if each message arrives while it is still fresh
\cite{simoens2018iort}. Surveys of networked embodied intelligence
organize its evaluation in levels, communication and computation
indicators below and embodied task indicators above, and argue that
only their joint measurement shows whether a network serves a task
\cite{zhang2026mean}; surveys of embodied multi-robot systems find that
benchmarks overlook communication latency and call for frameworks that
evaluate perception and communication across heterogeneous teams, with
hardware in the loop \cite{wen2026embodiedmrs}. Studying such a system
needs four things at
once: the network measured where the application runs, the channel
beneath it, every node's position and clock, and a way to put a new
application on the same network again. We review IoT testbeds, network
timeliness, wireless channel measurement, robots as IoT nodes, and
cooperative perception as the application, and close with the gap at
their intersection.

\subsection{IoT and Wireless Experimentation Testbeds}
Open testbeds made IoT and wireless experiments reproducible at scale.
FIT IoT-LAB \cite{adjih2015iotlab} reported 2,728 low-power nodes and
117 mobile robots across six indoor sites; its robots log their own
pose estimate, and its nodes record received power per channel. The
PAWR platforms POWDER \cite{breen2021powder} and COSMOS
\cite{raychaudhuri2020cosmos} provide city-scale software-defined
radio and edge computing, with small indoor labs beside the outdoor
deployment, and Colosseum \cite{bonati2021colosseum} emulates the RF
channel in FPGA hardware between 128 radios from modeled, ray-traced,
or measured scenarios. These testbeds make the radio the object of
study. Link metrics are instrumented by the experimenter, node
positions come from GNSS or the node itself, and no application task
is scored against ground truth, so they show what the network did but
not what it did to an application.

\subsection{Timeliness and Networking for IoT Applications}
The throughput of a contended 802.11 cell is modeled analytically
\cite{bianchi2000dcf}, and the freshness of status updates is captured
by the age of information \cite{kaul2012aoi,yates2021aoi}.
Task-oriented communication encodes features for a downstream
inference task rather than for reconstruction
\cite{shao2022taskoriented}, and edge computing for connected vehicles
\cite{liu2019edgeav} and cloud robotics \cite{kehoe2015cloud} place
the network on the critical path of perception. For robots, ROS~2
middleware (DDS, Zenoh, MQTT) has been benchmarked over Ethernet,
Wi-Fi, and 4G \cite{zhang2024ros2middleware}. Clocks are aligned by
NTP \cite{mills1991ntp} or, in sensor networks, by FTSP
\cite{maroti2004ftsp}. Measured networks are replayed by trace-driven
emulators \cite{noble1997trace,winstein2013sprout,netravali2015mahimahi,yan2018pantheon,hemminger2005netem}.
These works analyze, encode for, or emulate the network under
saturated, synthetic, or recorded traffic, but the deadline they
evaluate against is not that of a sensing application on the same
nodes, and a replayed trace feeds a socket, not a multi-node
application keyed to sensor time.

\subsection{Wireless Channel Measurement and Prediction}
Commodity CSI extraction, first on the Intel~5300
\cite{halperin2011csitool}, then with Nexmon \cite{nexmon} and
PicoScenes \cite{jiang2022picoscenes}, turned Wi-Fi into a sensor
\cite{ma2019wifisensing}. Channel datasets pair CSI or beam power with
positions or other modalities: DICHASUS \cite{dichasus} sounds a
distributed massive-MIMO channel indoors and outdoors with tachymeter
and GNSS positions, ESPARGOS \cite{espargos} records phase-coherent
Wi-Fi CSI from a robot-borne transmitter tracked by a total station,
MM-Fi \cite{mmfi} pairs CSI with human pose, and DeepSense~6G
\cite{deepsense6g} pairs mmWave beam power with camera, LiDAR, radar,
and RTK positions. Ray-traced digital twins
\cite{hoydis2023sionnart,alkhateeb2023digitaltwin,bostontwin,deepmimo}
and learned radio maps \cite{levie2021radiounet} predict the channel
from geometry. In this line, CSI is measured on dedicated sounding
links or used to sense people and passive targets, rather than on the
links over which IoT nodes exchange their application data; positions
come from a survey instrument or GNSS; and geometry is an input to
channel prediction, not an independent reference against which the
measured channel is checked.

\subsection{Robots as Mobile IoT Nodes}
The Internet of Robotic Things treats robots as mobile nodes that
sense, act, and share data over the same network as other devices
\cite{simoens2018iort}. Communication-aware robotics plans motion
around the channel \cite{muralidharan2021commaware}: robots map signal
strength online \cite{fink2010radiomap} or exploit its directionality
to keep a team connected \cite{gil2015adaptive}. CHORD and ACHORD
\cite{ginting2021chord,saboia2022achord} share data over intermittent
links in subterranean deployments of the DARPA SubT Challenge,
monitoring per-link SNR to route over dropped relays, and MOCHA
\cite{cladera2024mocha} does so in outdoor air--ground field trials,
synchronizing when the received signal strength allows. ROS-NetSim
\cite{calvofullana2021rosnetsim} routes real ROS traffic through a
simulated network. Multi-robot testbeds such as the Robotarium
\cite{wilson2021robotarium}, the CPM Lab \cite{kloock2021cpmlab}, and
ChoiRbot \cite{testa2021choirbot} obtain positions from overhead
cameras or motion capture and treat the network as ideal. Indoors
without motion capture, a pose reference needs surveyed UWB anchors
\cite{zhao2024util}, total-station points \cite{hilti2022}, surveyed
fiducials \cite{airmuseum}, or a scanner-built prior map
\cite{subtmrs,lamar}; SLAM-derived references are validated by sparse
landmark distances \cite{incrowdvi} or by internal consistency alone
\cite{envodat}. Robots therefore run over real links that are
monitored but not released, or have precise positions over a network
assumed ideal; none has both on one timeline, and none scores an
application against an independent reference.

\subsection{Cooperative Perception as an IoT Application}
Cooperative perception fuses raw points \cite{chen2019cooper},
features \cite{chen2019fcooper,wang2020v2vnet}, or selected regions
\cite{where2comm,liu2020when2com} from several nodes; F-Cooper
\cite{chen2019fcooper} frames it as an edge-computing system. Methods
now target what the network does to the fused result: latency
\cite{lei2022syncnet,v2xvit}, asynchrony \cite{wei2023cobevflow},
interruptions and packet loss \cite{ren2024v2xincop}, and pose error
\cite{lu2023coalign}. They are trained on simulated or outdoor data
\cite{opv2v,v2xsim,dairv2x} and evaluated against a constant delay, a
Bernoulli drop, or no latency, with pose error drawn from a chosen
distribution. Real-world platforms are outdoor and vehicular. V2V4Real
\cite{v2v4real} and TUMTraf-V2X \cite{zimmer2024tumtrafv2x} provide
GNSS-referenced multi-agent data with synchronization of tens of
milliseconds and no link log; \cite{hawlader2024v2xexp} reports V2X
experiments on real hardware; V2X-ReaLO \cite{v2xrealo} runs online
fusion over a live Wi-Fi link and reports its latency; and CooperScene
\cite{cooperscene}, the nearest case, releases pairwise C-V2X latency,
throughput, loss, and jitter traces with PTP-synchronized sensors, and
evaluates detectors under a delay derived from the measured mean
throughput. CoPeD \cite{coped}, an air--ground team recorded indoors
and outdoors, takes pose from its own visual-inertial odometry fused
with GPS and logs no link. The link, where recorded, is measured at or
above the transport layer without channel state, and it enters the
evaluation as a scalar delay or a live run, not as the per-message
trace each node actually had.

\subsection{Simulators and Environments for Embodied Agents}
Embodied agents are trained in simulators that generate their world:
Habitat \cite{habitat} for indoor navigation, CARLA \cite{carla} and
OpenCDA \cite{xu2021opencda} for driving, the latter generating the
cooperative-perception data of OPV2V \cite{opv2v}, and Isaac Lab
\cite{orbit} on Isaac Sim \cite{isaacsim} for robot learning at GPU
scale, behind the Gymnasium and PettingZoo interfaces
\cite{gymnasium,pettingzoo}. Lingjing \cite{lingjing} runs heterogeneous
aerial and ground agents in a simulated city behind a Gym-like
interface, with star or broadcast communication and episodes replayed
for failure attribution, but its messages pass an ideal bus whose cost
is a message count, never a delay or a loss. Driving research grounds
closed-loop simulation in real logs instead: nuPlan \cite{nuplan} and
Waymax \cite{waymax} replay recorded scenes around a simulated ego.
These environments either generate the world or replay it for one ego
over no network; none builds a multi-agent environment from a real
team's sensing and communication and checks it against that team.

\subsection{The Gap at the Intersection}
Table~\ref{tab:gap_iotj} compares the closest platforms on link,
channel, clock, pose reference, and replay; the supplementary material
states each gap by discipline, what SCooP-iT provides, and where the
paper shows it. IoT testbeds measure
the network but run no scored application; channel datasets measure
CSI on sounding links under a survey instrument; robot systems run
over real links but release no link log and score nothing; and
cooperative-perception platforms score the application outdoors under
GNSS but reduce the link to a scalar delay or a live run; and embodied
environments generate their world over an ideal network. SCooP-iT is
an indoor IoT testbed in which heterogeneous mobile and fixed nodes
record per-pair link and channel state on the sensors' clock, with a
measured bound on that clock, have a validated pose for every scan
without motion capture or GNSS, and replay the recorded per-message
link into any collaborative application, and build a closed-loop
environment from the same capture that is checked against it. It is the
integrated testbed the surveys call for
\cite{zhang2026mean,wen2026embodiedmrs}: sensing, communication, and
the task outcome of every run on one timeline, bridged to simulation by
an environment that answers to the real runs.

% Table II: related platforms, every cell checked on 18 Oct 2026 against
% official repos, project documentation, and excerpts of the papers
% (ACHORD paper, MOCHA and CPM Lab code read in full). Re-check against
% the PDFs before submission; arXiv and IEEE Xplore were not reachable.
\begin{table*}[!t]
\caption{Related IoT Testbeds, Channel Datasets, Networked Robots, and
Cooperative-Perception Platforms Compared on Link, Channel, Clock, Pose
Reference, and Replay}
\label{tab:gap_iotj}
\begin{center}
\renewcommand{\arraystretch}{1.15}
\setlength{\tabcolsep}{2pt}
\scriptsize
\begin{tabularx}{\textwidth}{@{}l l >{\raggedright\arraybackslash}p{1.3cm} c >{\raggedright\arraybackslash}X >{\raggedright\arraybackslash}p{1.1cm} >{\raggedright\arraybackslash}p{1.3cm} >{\raggedright\arraybackslash}p{1.4cm} >{\raggedright\arraybackslash}p{1.1cm} >{\raggedright\arraybackslash}p{1.1cm}@{}}
\toprule
\textbf{Platform} & \textbf{Venue} & \shortstack[l]{\textbf{Nodes}\\\textbf{(mob./fixed)}} & \textbf{Indoor} & \textbf{Measured link} & \textbf{CSI} & \shortstack[l]{\textbf{Pose}\\\textbf{reference}} & \shortstack[l]{\textbf{Clock}\\\textbf{sync}} & \shortstack[l]{\textbf{Link}\\\textbf{replay}} & \shortstack[l]{\textbf{App.}\\\textbf{scoring}} \\
\midrule
\multicolumn{10}{@{}l}{\textit{IoT and wireless testbeds}} \\
FIT IoT-LAB \cite{adjih2015iotlab} & WF-IoT'15 & 117 / 2728 & \cmark & Partial: RSSI per channel at node & \xmark & Self-estimate & Partial: GPS, subset & \xmark & \xmark \\
POWDER \cite{breen2021powder} & ComNet'21 & 20 / sites & partial & Partial: user-instrumented & Raw IQ & GNSS & Partial: GPSDO & \xmark & \xmark \\
COSMOS \cite{raychaudhuri2020cosmos} & MobiCom'20 & $\sim$200 / $\sim$50$^a$ & partial & Partial: user-instrumented & Raw IQ & Not reported & Not reported & \xmark & \xmark \\
Colosseum \cite{bonati2021colosseum} & DySPAN'21 & 128 emul. & --- & Emulated per pair & Emul.\ CIR & --- & Not reported & RF replay & \xmark \\
\midrule
\multicolumn{10}{@{}l}{\textit{Wi-Fi CSI tools and channel datasets}} \\
Nexmon CSI \cite{nexmon} & WiNTECH'19 & tool & --- & Partial: per-frame RSSI & \cmark & --- & \xmark & \xmark & \xmark \\
PicoScenes \cite{jiang2022picoscenes} & IoT-J'22 & tool & --- & Partial: per-frame RSSI, SNR & \cmark & --- & \xmark & \xmark & \xmark \\
DICHASUS \cite{dichasus} & WSA'21 & 1 / arrays & partial & --- & Sounding link & Tachymeter, DGNSS & Partial: OTA ref. & \xmark & \xmark \\
ESPARGOS \cite{espargos} & arXiv'24 & 1 / 4 arrays & \cmark & --- & Receive-only & Total station & Partial: shared clock & \xmark & \xmark \\
DeepSense 6G \cite{deepsense6g} & ComMag'23 & 1 / 1 & partial & Partial: beam power & \xmark & RTK & Partial: UTC stamps & \xmark & \xmark \\
\midrule
\multicolumn{10}{@{}l}{\textit{Robots as IoT nodes}} \\
Robotarium \cite{wilson2021robotarium} & RA-L'21 & $N$ / 0 & \cmark & --- & \xmark & Overhead tracking & \xmark & \xmark & \xmark \\
CPM Lab \cite{kloock2021cpmlab} & ECC'21 & $N$ / 0 & \cmark & Aggregate RTT (UI only) & \xmark & Overhead camera & NTP, no bound & \xmark & \xmark \\
ChoiRbot \cite{testa2021choirbot} & RA-L'21 & $N$ / 0 & \cmark & --- & \xmark & Motion capture & \xmark & \xmark & \xmark \\
CHORD / ACHORD \cite{ginting2021chord,saboia2022achord} & RA-L'21/'22 & 6 / base + relays & subterr. & Partial: per-link SNR, not released & \xmark & \xmark & \xmark & \xmark & \xmark \\
MOCHA \cite{cladera2024mocha} & ICRA'24 & UAV+UGVs / base & \xmark & Partial: per-peer RSSI, RTT, bandwidth, not released & \xmark & Not reported & \xmark & \xmark & \xmark \\
ROS-NetSim \cite{calvofullana2021rosnetsim} & RA-L'21 & sim. & --- & Simulated & \xmark & --- & \xmark & Sim. & \xmark \\
Mahimahi / NetEm \cite{netravali2015mahimahi,hemminger2005netem} & ATC'15 / LCA'05 & --- & --- & Recorded trace & \xmark & --- & \xmark & Transport & \xmark \\
\midrule
\multicolumn{10}{@{}l}{\textit{Cooperative-perception platforms}} \\
CoPeD \cite{coped} & RA-L'24 & 5 / 0 & partial & --- & \xmark & VIO + GPS & NTP, no number & \xmark & Partial$^b$ \\
V2V4Real \cite{v2v4real} & CVPR'23 & 2 / 0 & \xmark & --- & \xmark & GNSS/INS & Partial: $<$\,50\,ms & \xmark & Det., track. \\
TUMTraf-V2X \cite{zimmer2024tumtrafv2x} & CVPR'24 & 1 / 1 & \xmark & --- & \xmark & RTK & Partial: 24.9\,ms & \xmark & Det., track. \\
V2X-ReaLO \cite{v2xrealo} & arXiv'25 & 2 / 2 & \xmark & Latency (Wi-Fi) & \xmark & GNSS & No number & Live only & Det.\ (online) \\
CooperScene \cite{cooperscene} & ECCV'26 & 3 / 1 & \xmark & Pairwise C-V2X RTT, throughput, loss, jitter & \xmark & RTK + ICP & Claimed sub-ms & Scalar delay$^c$ & Detection \\
AutoCast \cite{qiu2022autocast} & MobiSys'22 & sim. & \xmark & Simulated & \xmark & Simulator & \xmark & Sim. & Sim. \\
\midrule
\textbf{SCooP-iT (ours)} & --- & 2 / 3 & \cmark & RSSI, MCS, retries (5\,Hz), RTT (100\,Hz), goodput (4\,s), per pair & \cmark$^d$ & Map + fiducials & Measured & Per message & C-SLAM, det. \\
\bottomrule
\end{tabularx}
\end{center}
\vspace{2pt}
{\scriptsize \cmark: provided; partial: provided in part, or only if the experimenter instruments it; \xmark: not provided; ---: not applicable.\quad
$^a$Planned deployment in the 2020 paper.\quad
$^b$Zero-shot foundation-model image labels, no 3D boxes.\quad
$^c$One delay derived from mean measured throughput, not a per-message trace.\quad
$^d$Single-antenna Wi-Fi CSI on the data links (BCM43455c0, \SI{80}{\mega\hertz}, \SIrange{104}{141}{\hertz}), not phase-coherent across nodes.}
\end{table*}

